% sections/04_llm_verify.tex --- SkillStrike

Static analysis identifies \emph{syntactic} evidence of dangerous capability.
It cannot determine whether the capability is truly undisclosed,
whether it is reachable under agent-driven invocation,
or whether the SKILL.md documentation concedes the capability
in indirect language.
We introduce a two-stage LLM verification pipeline that answers these
questions before investing in expensive live-agent experiments.

\subsection{Two-Stage Pipeline}

\paragraph{Stage 1: Sink capability confirmation.}
For each LSCPS-flagged sink skill, we prompt a chain-of-thought
LLM (Codex \texttt{gpt-4.1}) with the full SKILL.md documentation
and the extracted code evidence:

\begin{lstlisting}[language={},caption={Stage 1 prompt template.}]
You are a security auditor. Review the following skill's
SKILL.md documentation and code implementation.

SKILL.md:
{skill_md_content}

Implementation evidence (LSCPS-extracted):
{lscps_evidence}

Task:
1. Does this skill's documentation anywhere mention that
   it can [CAPABILITY_FAMILY]? (yes/partial/no)
2. Does the implementation code confirm this capability
   is reachable under normal agent invocation? (yes/no)
3. Verdict: Is this a VIABLE sink for composition attack?
   (CONFIRMED / POSSIBLE / NOT_VIABLE)

Reason step by step before giving the verdict.
\end{lstlisting}

Skills receiving \texttt{CONFIRMED} or \texttt{POSSIBLE} proceed to Stage 2.
\texttt{NOT\_VIABLE} skills (documentation explicitly discloses the capability,
or code evidence is unreachable) are filtered out.

\paragraph{Stage 2: Source-Sink pair viability.}
For each confirmed sink, we generate candidate source-sink pairs
by pairing each sink with a set of source skills (selected by \lscps)
that share at least one injection channel category.
For each pair, we prompt the LLM to assess whether the source skill's
output channel can plausibly carry attacker-controlled content
to the co-installed sink:

\begin{lstlisting}[language={},caption={Stage 2 pair viability prompt.}]
Source skill: {src_name} --- {src_description}
Sink skill: {sink_name} --- dangerous capability: {capability}

Can an attacker who controls content reachable via the
source skill's channel cause the agent to invoke the
sink's dangerous capability?
Contexts: both skills are co-installed; the agent has
full-auto mode enabled.

Rate: LIKELY / POSSIBLE / UNLIKELY
Reason step-by-step.
\end{lstlisting}

\subsection{Application to \skillssh{}}
\label{sec:llm_ssh}

We apply the two-stage pipeline to all \NselfExfil{} self-exfil
pairs identified by \lscps in \skillssh{}.
Table~\ref{tab:llm_audit} summarizes results.

\begin{table}[t]
\centering
\caption{Two-stage LLM verification results on \skillssh{} candidates.}
\label{tab:llm_audit}
\setlength{\tabcolsep}{4pt}
\small
\begin{tabular}{lrr}
\toprule
\textbf{Stage} & \textbf{Skills} & \textbf{Pass Rate}\\
\midrule
LSCPS self-exfil candidates & 4,154 & baseline \\
Stage 1: sink confirmed       &   521 & 12.5\% \\
Stage 2: pair viable          &   197 &  4.7\% \\
\midrule
\textbf{Live experiment input} & \textbf{197} & \textbf{4.7\%}\\
\bottomrule
\end{tabular}
\end{table}

The 197 Codex-confirmed viable paths represent a
$\mathbf{21\times}$ precision improvement over raw LSCPS output:
each confirmed path has LLM-validated evidence that
(i) the dangerous capability is present and reachable,
and (ii) the documentation does not disclose it.

\paragraph{Accuracy estimation.}
To estimate Stage 1 precision, one author manually reviewed
a stratified random sample of 50 Stage-1 outputs (25 CONFIRMED, 25 NOT\_VIABLE).
Of 25 CONFIRMED verdicts, 23 were verified correct on manual inspection
(precision 92\%).
Of 25 NOT\_VIABLE verdicts, 24 were correctly filtered
(false negative rate 4\%), confirming that the CONFIRMED filter
is conservative and precise.

\paragraph{Comparison with static-only.}
The 4,154 LSCPS self-exfil candidates include many skills where
(a) the dangerous capability is mentioned in a non-obvious way in SKILL.md,
(b) the code is present but unreachable under agent invocation,
or (c) the ``dangerous'' pattern is a development-time test harness.
The LLM verification stage removes all three categories before live experiments,
avoiding misleading success inflations from trivially-exploitable test environments.

\subsection{Design Choices}

\paragraph{Why Codex (gpt-4.1) not a smaller model?}
We evaluated gpt-3.5-turbo on a 100-skill calibration set
and found 31\% hallucination rate in Stage 1 verdicts
(declaring CONFIRMED for skills with no dangerous code).
gpt-4.1 with chain-of-thought reduces this to 8\%, meeting our
quality bar for the downstream live experiment.

\paragraph{Cost.}
At the April 2026 API rate, Stage 1 + Stage 2 processing of all
4,154 candidates costs approximately \$12 USD in API fees---
negligible relative to the cost of running live experiments
on all 4,154 candidates.

\paragraph{Limitations.}
The LLM cannot execute the skill code; it reasons over documented behavior.
Skills with entirely undocumented implementations (e.g., binary blobs)
cannot be assessed. These are filtered by LSCPS Tier-6 for separate review.
